% Part: sets-functions-relations
% Chapter: arithmetization
% Section: rationals

\documentclass[../../../include/open-logic-section]{subfiles}

\begin{document}

\olfileid[tr]{sfr}{arith}{rat}
\olsection{$\Int$'den $\Rat$'ye}

Az önce biraz naif küme teorisi kullanarak tam sayıların doğal sayılardan nasıl
kurulacağını gördük. Şimdi de çok benzer bir yolla rasyonel sayıların tam
sayılardan nasıl kurulacağını göreceğiz. Başlangıç gözlemlerimiz şunlardır:
\begin{enumerate}
	\item Her rasyonel sayı, $i$ ve $j$ tam sayılar, $j$ ise sıfırdan farklı
	olmak üzere $\nicefrac{i}{j}$ biçiminde yazılabilir.
	\item $\nicefrac{i}{j}$ ifadesinde kodlanan bilgi, $\tuple{ i, j}$ sıralı
	ikilisiyle de kodlanabilir.
\end{enumerate}
Açık yaklaşım, rasyonel sayıları $\Int \times (\Int \setminus \{0_\Int\})$'den
alınan sıralı ikililer \emph{olarak} düşünmek olurdu. Ancak daha önce olduğu
gibi bu da biraz fazla naif olurdu; çünkü $\nicefrac{3}{2} =
\nicefrac{6}{4}$ olmasını isteriz, fakat $\tuple{ 3, 2}\neq \tuple{ 6, 4}$.
Daha genel olarak şunu isteriz:
\[
	\nicefrac{a}{b} = \nicefrac{c}{d} \text{ ancak ve ancak } a \times d = b \times c
\]
Bunu elde etmek için $\Int \times (\Int \setminus \{0_\Int\})$ üzerinde
şöyle bir {denklik bağıntısı} tanımlarız:
\[
	\tuple{ a, b }\Ratequiv \tuple{ c, d} \text{ ancak ve ancak }a \times d = b \times c
\]
Bunun bir denklik bağıntısı olduğunu denetlemeliyiz. Bu denetim
$\Intequiv$ durumuna çok benzer; onu alıştırma olarak bırakacağız.
\begin{prob}
$\Ratequiv$ bağıntısının bir denklik bağıntısı olduğunu gösterin.
\end{prob}
Bu bağıntı sayesinde şunu söyleyebiliriz:
\begin{defn}
Rasyonel sayılar, ikinci bileşeni sıfırdan farklı olan tam sayı ikililerinin
$\Ratequiv$ altındaki denklik sınıflarıdır. Yani
$\Rat = \equivclass{(\Int \times (\Int\setminus \{0_\Int\}))}{\Ratequiv}$.
\end{defn}

Tam sayılarda olduğu gibi bazı temel işlemleri de tanımlamak isteriz.
$\equivrep{i,j}{\Ratequiv}$, $\tuple{i, j}$'yi !!a{element} olarak içeren
$\Ratequiv$ denklik sınıfı olmak üzere şöyle deriz:
\begin{align*}
	\equivrep{a, b}{\Ratequiv} + \equivrep{c, d}{\Ratequiv} &= \equivrep{ad + bc,  bd}{\Ratequiv}\\
	\equivrep{a, b}{\Ratequiv} \times \equivrep{c, d}{\Ratequiv} &= \equivrep{a  c, b d}{\Ratequiv}.
\intertext{Bu rasyonel sayılar üzerinde $r \leq s$ bağıntısını tanımlamak için
$r \le s$'nin ancak ve ancak $s-r$ negatif değilse geçerli olduğu; başka bir
deyişle $s-r$'nin, $i$ negatif olmayan ve $j$ pozitif olan
$\nicefrac{i}{j}$ biçiminde yazılabileceği olgusunu kullanırız:}
	\equivrep{a, b}{\Ratequiv} \leq \equivrep{c, d}{\Ratequiv} &\text{ ancak ve ancak }
	\equivrep{c, d}{\Ratequiv} - \equivrep{a, b}{\Ratequiv} =
	\equivrep{i_\Int, j_\Int}{\Ratequiv}
\end{align*}
burada bazı $i \in \Nat$ ve $0 \neq j \in \Nat$ vardır.

Sonra bu tanımların \emph{gerektiği gibi} davrandığını denetlememiz gerekir;
bu denetimi \olref[check]{sec}'e bırakıyoruz. Ama gerçekten de öyle
davranırlar!{} Son olarak tam sayıları rasyonel sayılar \emph{olarak} ele
almanın bir yolunu isteriz; bu nedenle her $i \in \Int$ için
$i_\Rat = \equivrep{i, 1_\Int}{\Ratequiv}$ olduğunu belirleriz. Yine bütün
bunların doğru davrandığını \olref[check]{sec}'te denetliyoruz.

\begin{prob}
Her $i, j \in \Int$ için $(i + j)_\Rat = i_\Rat+ j_\Rat$,
$(i \times j)_\Rat = i_\Rat \times j_\Rat$ ve
$i \leq j \liff i_\Rat \leq j_\Rat$ olduğunu gösterin.
\end{prob}

\end{document}
